% Options for packages loaded elsewhere
\PassOptionsToPackage{unicode}{hyperref}
\PassOptionsToPackage{hyphens}{url}
\documentclass[
  ignorenonframetext,
]{beamer}
\newif\ifbibliography
\usepackage{pgfpages}
\setbeamertemplate{caption}[numbered]
\setbeamertemplate{caption label separator}{: }
\setbeamercolor{caption name}{fg=normal text.fg}
\beamertemplatenavigationsymbolsempty
% remove section numbering
\setbeamertemplate{part page}{
  \centering
  \begin{beamercolorbox}[sep=16pt,center]{part title}
    \usebeamerfont{part title}\insertpart\par
  \end{beamercolorbox}
}
\setbeamertemplate{section page}{
  \centering
  \begin{beamercolorbox}[sep=12pt,center]{section title}
    \usebeamerfont{section title}\insertsection\par
  \end{beamercolorbox}
}
\setbeamertemplate{subsection page}{
  \centering
  \begin{beamercolorbox}[sep=8pt,center]{subsection title}
    \usebeamerfont{subsection title}\insertsubsection\par
  \end{beamercolorbox}
}
% Prevent slide breaks in the middle of a paragraph
\widowpenalties 1 10000
\raggedbottom
\AtBeginPart{
  \frame{\partpage}
}
\AtBeginSection{
  \ifbibliography
  \else
    \frame{\sectionpage}
  \fi
}
\AtBeginSubsection{
  \frame{\subsectionpage}
}
\usepackage{iftex}
\ifPDFTeX
  \usepackage[T1]{fontenc}
  \usepackage[utf8]{inputenc}
  \usepackage{textcomp} % provide euro and other symbols
\else % if luatex or xetex
  \usepackage{unicode-math} % this also loads fontspec
  \defaultfontfeatures{Scale=MatchLowercase}
  \defaultfontfeatures[\rmfamily]{Ligatures=TeX,Scale=1}
\fi
\usepackage{lmodern}
\ifPDFTeX\else
  % xetex/luatex font selection
\fi
% Use upquote if available, for straight quotes in verbatim environments
\IfFileExists{upquote.sty}{\usepackage{upquote}}{}
\IfFileExists{microtype.sty}{% use microtype if available
  \usepackage[]{microtype}
  \UseMicrotypeSet[protrusion]{basicmath} % disable protrusion for tt fonts
}{}
\makeatletter
\@ifundefined{KOMAClassName}{% if non-KOMA class
  \IfFileExists{parskip.sty}{%
    \usepackage{parskip}
  }{% else
    \setlength{\parindent}{0pt}
    \setlength{\parskip}{6pt plus 2pt minus 1pt}}
}{% if KOMA class
  \KOMAoptions{parskip=half}}
\makeatother
\setlength{\emergencystretch}{3em} % prevent overfull lines
\providecommand{\tightlist}{%
  \setlength{\itemsep}{0pt}\setlength{\parskip}{0pt}}
\usepackage{bookmark}
\IfFileExists{xurl.sty}{\usepackage{xurl}}{} % add URL line breaks if available
\urlstyle{same}
\hypersetup{
  hidelinks,
  pdfcreator={LaTeX via pandoc}}

\author{\texorpdfstring{}{}}
\date{}

\begin{document}

\begin{frame}[fragile]{Methodology}
\protect\phantomsection\label{sec:methodology}
This paper presents a theoretical framework paper rather than an
empirical study. The methodology combines conceptual analysis,
mathematical formalization, and computational validation to establish
and evaluate the \(2 \times 2\) quadrant structure for Active
Inference's meta-level operation.

\begin{block}{Theoretical Approach}
\protect\phantomsection\label{theoretical-approach}
The research adopts a conceptual-analytic methodology grounded in the
Active Inference literature \cite{friston2010free, parr2022active}. The
\(2 \times 2\) framework was derived by identifying two orthogonal
dimensions along which Active Inference's contributions can be
organized:

\textbf{Data/Meta-Data dimension.} This axis reflects the information
hierarchy: \emph{Data} refers to raw sensory observations processed
through the generative model's likelihood mapping (\(A\) matrix), while
\emph{Meta-Data} refers to higher-order information about data quality,
reliability, and provenance (confidence scores, temporal stamps, source
credibility). The distinction captures the difference between what is
observed and what is known about the observation process itself.

\textbf{Cognitive/Meta-Cognitive dimension.} This axis reflects the
processing level hierarchy, following the meta-cognitive tradition from
\cite{flavell1979metacognition} and \cite{nelson1990metamemory}:
\emph{Cognitive} processing involves direct transformation of
information through belief updating and policy selection, while
\emph{Meta-Cognitive} processing involves monitoring, evaluating, and
regulating cognitive processes themselves. In Active Inference terms,
this distinguishes EFE computation from reasoning about EFE computation.

The cross-product of these dimensions yields four quadrants (Q1--Q4),
each representing a distinct mode of cognitive operation with specific
mathematical formulations grounded in the EFE framework.
\end{block}

\begin{block}{Mathematical Formalization}
\protect\phantomsection\label{mathematical-formalization}
Each quadrant is formalized through extensions of the standard EFE
formulation (Equation \eqref{eq:efe}):

\begin{itemize}
\tightlist
\item
  \textbf{Q1} employs the baseline EFE decomposition into epistemic and
  pragmatic components.
\item
  \textbf{Q2} extends EFE with meta-data weighting parameters (\(w_e\),
  \(w_p\), \(w_m\)) that modulate the relative influence of epistemic,
  pragmatic, and meta-data-derived value.
\item
  \textbf{Q3} introduces hierarchical EFE with a meta-cognitive term
  \(\mathcal{F}_{meta}(\pi)\) weighted by \(\lambda\), enabling
  self-assessment and adaptive strategy selection.
\item
  \textbf{Q4} formulates framework-level optimization over the parameter
  space \(\Theta\), with regularization \(\mathcal{R}(\Theta)\) ensuring
  coherence.
\end{itemize}

All formulations maintain consistency with the variational free energy
framework and reduce to standard Active Inference in appropriate limits.
\end{block}

\begin{block}{Computational Validation}
\protect\phantomsection\label{computational-validation}
The theoretical framework is validated through computational
implementations in the project's source modules. Each quadrant is
implemented as a testable module with deterministic outputs (fixed
random seeds) enabling reproducible validation:

\begin{itemize}
\tightlist
\item
  \textbf{Quadrant implementations}
  (\texttt{src/framework/quadrant\_framework.py}) instantiate each
  quadrant's mathematical formulation with concrete matrix
  specifications.
\item
  \textbf{Active inference core}
  (\texttt{src/core/active\_inference.py}) implements EFE calculation,
  policy selection, and belief updating.
\item
  \textbf{Meta-cognitive module}
  (\texttt{src/framework/meta\_cognition.py}) implements confidence
  assessment, adaptive attention allocation, and strategy selection.
\item
  \textbf{Generative model specification}
  (\texttt{src/core/generative\_models.py}) implements \(A\), \(B\),
  \(C\), \(D\) matrix construction and validation.
\item
  \textbf{Visualization engine}
  (\texttt{src/visualization/visualization.py}) produces
  publication-quality figures with programmatic font and layout
  enforcement.
\end{itemize}

Validation metrics include decision accuracy under uncertainty,
confidence calibration quality, strategy adaptation appropriateness, and
framework optimization convergence. Statistical validation employs
hypothesis testing (t-tests, ANOVA) with effect sizes reported as
Cohen's \(d\), as detailed in Appendix \ref{sec:validation_results}.
\end{block}

\begin{block}{Scope and Limitations}
\protect\phantomsection\label{scope-and-limitations}
As a theoretical framework paper, the methodology is subject to inherent
limitations: the quadrant boundaries are analytical constructs that may
not correspond to sharp neural or computational distinctions, and the
computational validation demonstrates feasibility rather than empirical
necessity. We address these limitations explicitly in Section
\ref{sec:limitations}.
\end{block}
\end{frame}

\end{document}
